% glossary.tex -- Module 5 ESSENTIALS glossary (A-Z): the key terms only.
% src/mod5/glossary.tex (term, meaning and Chinese copied word for word, except the
% rows corrected in the script's FIXES). Edit the script, not this file.
\unnumberedchapter[\faBook]{Glossary}

The key words of this module, from A to Z. The meaning is in plain English. The last column
gives the usual Chinese term used in physics textbooks. The full guide's glossary has more
words.

\begin{glossarytable}
  % ---- A
  \gterm{Acceleration due to gravity ($g$)}{the acceleration of a falling object near the Earth: \qty{9.8}{m.s^{-2}} downwards}{重力加速度}
  \gterm{Angular displacement ($\Delta\theta$)}{the angle an object turns through, measured at the centre}{角位移}
  \gterm{Angular velocity ($\omega$)}{the angle turned each second, in \si{rad.s^{-1}}}{角速度}
  \gterm{Aphelion}{the farthest point of an orbit from the Sun}{远日点}
  \gterm{Apogee}{the farthest point of an orbit around the Earth}{远地点}
  \gterm{Apparent weightlessness}{the feeling of having no weight in free fall, because the normal force is zero}{完全失重}
  % ---- B
  \gterm{Balanced torques}{clockwise and anticlockwise torques are equal, so the rotation does not change}{力矩平衡}
  \gterm{Banked track}{a curved road that is tilted, higher on the outside of the bend}{倾斜弯道（外高内低）}
  \gterm{Bound}{trapped by gravity, so the object cannot get away; its total energy is negative}{被引力束缚的}
  % ---- C
  \gterm{Centrifugal force}{a made-up ``outward force''; do not use it: the outward feeling is inertia}{离心力}
  \gterm{Centripetal acceleration ($a_\mathrm{c}$)}{the acceleration towards the centre of a circle: $v^2/r$}{向心加速度}
  \gterm{Centripetal force ($F_\mathrm{c}$)}{the net force towards the centre that keeps an object moving in a circle}{向心力}
  \gterm{Coefficient of friction ($\mu$)}{a number with no unit that says how strongly two surfaces grip each other}{摩擦因数}
  \gterm{Component}{the part of a vector (such as a force) that points in one chosen direction}{分量（分力）}
  \gterm{Conical pendulum}{a mass on a string moving in a horizontal circle; the string sweeps out a cone}{圆锥摆}
  % ---- E
  \gterm{Ellipse}{a stretched circle (an oval)}{椭圆}
  \gterm{Escape velocity ($v_\mathrm{esc}$)}{the smallest launch speed that lets a projectile get away from a planet for ever}{逃逸速度}
  % ---- F
  \gterm{Free fall}{motion when gravity is the only force acting}{自由下落（只受引力）}
  % ---- G
  \gterm{Geostationary orbit}{a circular orbit above the equator with a one-day period, so the satellite stays above one point}{地球静止轨道}
  \gterm{Geosynchronous orbit}{any orbit with a period of one day; it may be tilted}{地球同步轨道}
  \gterm{Gravitational field strength ($g$)}{the gravitational force on each kilogram at a point, in \si{N.kg^{-1}} (\si{m.s^{-2}})}{引力场强度}
  \gterm{Gravitational force}{the pull between any two objects that have mass}{引力（万有引力）}
  \gterm{Gravitational potential energy ($U$)}{energy stored because of where an object is in a gravitational field; near the ground it changes by $mg\Delta h$; in space $U = -GMm/r$}{引力势能（地面附近：重力势能）}
  % ---- H
  \gterm{Hohmann transfer orbit}{a fuel-saving half-ellipse between two circular orbits, using two short engine burns (both forwards to move up, both backwards to move down)}{霍曼转移轨道}
  % ---- I
  \gterm{Inertia}{the tendency of an object to stay at rest, or to keep moving in a straight line at a steady speed}{惯性}
  \gterm{Infinity}{a place so far away that the planet's gravity there is too small to matter}{无穷远处}
  \gterm{Inverse-square law}{a rule where doubling the distance makes the effect four times smaller}{平方反比定律}
  % ---- K
  \gterm{Kepler's first law (law of ellipses)}{each planet moves in an ellipse with the Sun at one focus}{开普勒第一定律（轨道定律）}
  \gterm{Kepler's second law (law of equal areas)}{the line from the Sun to a planet sweeps out equal areas in equal times}{开普勒第二定律（面积定律）}
  \gterm{Kepler's third law (law of periods)}{$T^2$ is proportional to $r^3$ for everything that orbits the same body}{开普勒第三定律（周期定律）}
  \gterm{Kinetic energy ($K$)}{energy of motion: $K = \tfrac{1}{2}mv^2$}{动能}
  % ---- L
  \gterm{Line of action}{the line along which a force acts, drawn longer in both directions}{力的作用线}
  \gterm{Low Earth orbit (LEO)}{an orbit about \qtyrange{200}{2000}{km} above the Earth, with a period of about 90 minutes to 2 hours}{低地球轨道（近地轨道）}
  % ---- M
  \gterm{Maximum height ($h_\mathrm{max}$)}{the greatest height above the launch point}{最大高度}
  % ---- N
  \gterm{Net force}{the total of all the forces on an object, added with their directions}{合力}
  \gterm{Normal force ($F_\mathrm{N}$)}{the push of a surface on an object, at \ang{90} to the surface}{支持力（法向力）}
  % ---- O
  \gterm{Orbital decay}{the slow fall to lower orbits caused by drag from the thin upper atmosphere}{轨道衰减}
  \gterm{Orbital speed ($v$)}{the speed a satellite needs to stay in a circular orbit}{环绕速度（轨道速度）}
  % ---- P
  \gterm{Perigee}{the closest point of an orbit around the Earth}{近地点}
  \gterm{Perihelion}{the closest point of an orbit to the Sun}{近日点}
  \gterm{Period ($T$)}{the time for one full turn}{周期}
  \gterm{Perpendicular distance ($r_\perp$)}{the shortest distance from the pivot to the line of action of a force (the lever arm)}{力臂}
  \gterm{Pivot}{the fixed point or line that an object turns about, or that a string hangs from}{支点（转轴；悬点）}
  \gterm{Projectile}{an object that is thrown or launched and then moves with no engine; in the model, gravity is the only force on it}{抛体}
  % ---- R
  \gterm{Radial field}{a field whose lines all point straight at the centre of the planet, like the spokes of a wheel}{径向场}
  \gterm{Radian (rad)}{a unit for angles: one full turn is $2\pi$ rad, the same as \ang{360}}{弧度}
  \gterm{Range ($R$)}{the horizontal distance from the launch point to the landing point}{射程（水平射程）}
  % ---- S
  \gterm{Semi-major axis ($a$)}{half of the longest diameter of an ellipse}{半长轴}
  \gterm{Static friction ($F_\mathrm{f}$)}{friction that stops surfaces sliding; it grows as needed, up to a maximum of $\mu F_\mathrm{N}$}{静摩擦力}
  % ---- T
  \gterm{Tangent}{a line that touches a circle at one point, at \ang{90} to the radius}{切线}
  \gterm{Tension ($F_\mathrm{T}$)}{the pull of a string, rope or chain, along its length}{拉力（张力）}
  \gterm{Time of flight ($t_\mathrm{flight}$)}{the time from launch until landing}{飞行时间}
  \gterm{Torque ($\tau$)}{the turning effect of a force, in \si{N.m}}{力矩}
  \gterm{Total mechanical energy ($E$)}{kinetic energy plus potential energy: $E = K + U$}{机械能}
  % ---- U
  \gterm{Uniform circular motion}{moving around a circle at a constant speed}{匀速圆周运动}
  \gterm{Uniform velocity}{a velocity that does not change: the same speed in the same direction; the acceleration is zero}{匀速直线运动（速度恒定）}
  \gterm{Universal gravitational constant ($G$)}{the fixed number that sets how strong gravity is, \qty{6.67e-11}{N.m^2.kg^{-2}}}{引力常量}
  % ---- W
  \gterm{Work ($W$)}{energy transferred by a force when the object moves along the direction of the force (or part of it)}{功}
  \gterm{Work--energy theorem}{the work done by the net force equals the change in kinetic energy}{动能定理}
\end{glossarytable}
